\documentclass[12pt,a4paper]{article}

% Packages
\usepackage[utf8]{inputenc}
\usepackage[T1]{fontenc}
\usepackage{amsmath,amssymb,amsfonts}
\usepackage{mathtools}
\usepackage{graphicx}
\usepackage{booktabs}
\usepackage{hyperref}
\usepackage[margin=2.5cm]{geometry}
\usepackage{natbib}
\usepackage{algorithm}
\usepackage{algpseudocode}
\usepackage{xcolor}
\usepackage{listings}
\usepackage{subcaption}
\usepackage{siunitx}
\usepackage{multirow}

% Listing style
\lstset{
  basicstyle=\footnotesize\ttfamily,
  keywordstyle=\color{blue},
  commentstyle=\color{gray},
  breaklines=true,
  frame=single,
  numbers=left,
  numberstyle=\tiny\color{gray}
}

% Custom commands
\newcommand{\Neff}{N_{\mathrm{eff}}}
\newcommand{\Nreal}{N_{\mathrm{real}}}
\DeclareMathOperator{\MI}{MI}

\title{Numerical Validation of a Pure Python Implementation of k-NN Entropy Estimation Against a Reference C/ANN Library}
\author{Nicolas B.\ Garnier$^{1}$ \and Martin G.\ Frasch$^{2}$}

\newcommand{\affiliations}{%
$^{1}$Laboratoire de Physique, CNRS, ENS de Lyon, Universit\'{e} de Lyon, Lyon, France\\
$^{2}$Department of Obstetrics and Gynecology, University of Washington, Seattle, WA, USA
}
\date{\today}

\begin{document}
\maketitle
\begingroup
\centering
\small
\affiliations
\par
\endgroup
\vspace{1em}

\begin{abstract}
We present a systematic validation of a pure Python implementation of k-nearest-neighbor (k-NN) based entropy and entropy-rate estimators against a mature reference implementation in C that uses the Approximate Nearest Neighbor (ANN) library.
Testing on 251 segments of RR-interval time series from ambulatory ECG recordings, we identify and resolve three sources of numerical discrepancy: (i)~the ANN library's exclusion of all zero-distance points from neighbor searches, (ii)~a mismatch between $L_2$ and $L_\infty$ norms in high-dimensional k-d tree queries, and (iii)~a misalignment of data points between the marginal-entropy and mutual-information terms in the entropy-rate computation at time scales $\tau > 1$.
After correction, Shannon entropy estimates match to machine precision ($R^2 = 1.0000$) and entropy-rate estimates agree to within $0.06\%$ ($R^2 = 0.9998$) across all tested time scales, confirming numerical equivalence of the two codebases.
These findings highlight subtle implementation pitfalls in k-NN information-theoretic estimators that are relevant to any reimplementation effort.
\end{abstract}

\section{Introduction}\label{sec:intro}

Entropy and entropy-rate estimation from finite samples is central to nonlinear time-series analysis in physiology, neuroscience, and engineering~\citep{kozachenko1987,kraskov2004}.
The Kozachenko--Leonenko (KL) estimator~\citep{kozachenko1987} and the Kraskov--St\"ogbauer--Grassberger (KSG) mutual information estimator~\citep{kraskov2004} provide consistent, bias-corrected estimates based on k-nearest-neighbor (k-NN) distances.
While conceptually straightforward, their numerical implementation involves several choices---distance metric, treatment of zero-distance neighbors, time-delay embedding alignment---that can introduce systematic offsets when different codebases are compared.

We describe the validation of a new pure Python implementation (\texttt{entropy\_pure}) against an established C implementation that relies on the ANN library~\citep{mount2010} for k-d tree operations.
The reference C code has been used extensively in published studies of heart-rate variability (HRV) and fetal monitoring~\citep{mayer2025}.
The Python reimplementation aims to provide a dependency-free, readable, and extensible alternative suitable for research prototyping and teaching, while producing numerically identical results.

The validation uncovered three implementation-level discrepancies that collectively produced offsets of up to 14\% in entropy-rate estimates.
We document these discrepancies, their root causes, and the corrections applied, as a practical guide for developers of information-theoretic estimation libraries.


\section{Methods}\label{sec:methods}

\subsection{Kozachenko--Leonenko entropy estimator}\label{sec:kl}

Given $N$ i.i.d.\ samples $\{x_1, \ldots, x_N\}$ in $\mathbb{R}^d$, the KL estimator of the differential entropy $H(X)$ is
\begin{equation}\label{eq:kl}
  \hat{H}(X) = d \cdot \frac{1}{N_{\mathrm{v}}} \sum_{\substack{i=1 \\ \varepsilon_i > 0}}^{N} \log(2\varepsilon_i)
  + \psi(N_{\mathrm{v}}) - \psi(k),
\end{equation}
where $\varepsilon_i$ is the distance from $x_i$ to its $k$-th nearest neighbor, $\psi(\cdot)$ is the digamma function, and $N_{\mathrm{v}} \leq N$ is the number of points with $\varepsilon_i > 0$.
The factor of~2 converts the distance to the diameter of the $L_\infty$ ball, and the digamma terms provide the bias correction.

Two choices in this estimator have significant numerical impact:
\begin{enumerate}
  \item \textbf{Distance metric.} The C implementation uses the ANN library compiled with $L_\infty$ (Chebyshev) norm (\texttt{ANN\_METRIC = LINF}).
  While $L_2 = L_\infty$ in one dimension, the two norms diverge for $d \geq 2$, producing a systematic offset in the Shannon entropy $H$ at embedding dimension $m \geq 2$.
  \item \textbf{Zero-distance handling.} When \texttt{ANN\_ALLOW\_SELF\_MATCH = false}, the ANN k-d tree search excludes \emph{all} points at distance zero from the query result---not merely the query point itself.
  For integer-valued data with many duplicate values (common in RR-interval time series quantized to \SI{1}{\milli\second}), this means the $k$-th neighbor always has $\varepsilon_i > 0$, whereas a standard k-d tree (e.g., \texttt{scipy.spatial.KDTree}) returns exact zeros for duplicate points.
\end{enumerate}

\subsection{KSG mutual information estimator}\label{sec:ksg}

For the entropy-rate computation via method~2 ($h = H(X_t) - \MI(X_t; X_{\mathrm{past}})$), we use the KSG estimator (Algorithm~1):
\begin{equation}\label{eq:ksg}
  \hat{I}(X; Y) = \psi(k) - \frac{1}{N} \sum_{i=1}^{N} \bigl[\psi(n_{x,i} + 1) + \psi(n_{y,i} + 1)\bigr] + \psi(N),
\end{equation}
where $n_{x,i}$ and $n_{y,i}$ count the number of points within the $L_\infty$-ball of radius $\varepsilon_i$ in each marginal space, with $\varepsilon_i$ the $L_\infty$ distance to the $k$-th neighbor in the joint space.

\subsection{Entropy rate at multiple time scales}\label{sec:erate}

The entropy rate at time scale $\tau$ is
\begin{equation}\label{eq:erate}
  \hat{h}(\tau) = \hat{H}(X_t) - \hat{I}(X_t;\, X_{\mathrm{past}}^{(m)}),
\end{equation}
where $X_{\mathrm{past}}^{(m)} = (X_{t-\tau}, X_{t-2\tau}, \ldots, X_{t-m\tau})$ is the $m$-dimensional past embedding with lag~$\tau$.

The reference C implementation constructs a \emph{unified embedding} of $m+1$ dimensions via its \texttt{Theiler\_embed} function:
for each selected time index $t$, the embedded vector is $(x_t,\, x_{t-\tau},\, \ldots,\, x_{t-m\tau})$.
Crucially, consecutive embedded vectors step by one time unit (the Theiler window $\Delta_T = 1$), not by $\tau$.
Both $\hat{H}(X_t)$ and $\hat{I}(X_t; X_{\mathrm{past}})$ operate on the \emph{same} set of $\Neff$ points extracted from this unified embedding, ensuring algebraic consistency of the subtraction in Eq.~\eqref{eq:erate}.

\subsection{Data and experimental setup}\label{sec:data}

We use 251 five-minute segments of RR-interval time series extracted from ambulatory fetal ECG recordings.
Segment lengths range from 139 to 547 beats (mean 422).
The RR intervals are integer-valued in milliseconds, with substantial duplicate values: a typical 200-point subsample contains approximately 120 unique values, causing ${\sim}40\%$ of $k$-th neighbor distances to be exactly zero when using a standard k-d tree.

All comparisons use $k = 5$, $\Neff = 200$, $N_{\mathrm{real}} = 1$, embedding dimension $m = 2$, and explicit Theiler window $\Delta_T = 1$ (``type~0'' in the C code's sampling-parameter system).
Time scales tested are $\tau \in \{1, 2, 4\}$.


\section{Results}\label{sec:results}

\subsection{Discrepancy 1: Zero-distance duplicate exclusion}\label{sec:bug1}

\subsubsection{Observation}
The one-dimensional marginal entropy $H(X_t)$ at $\tau = 1$ showed a mean offset of \SI{0.43}{nats} between the Python and C implementations ($R^2 = 0.836$).
Decomposing the entropy rate $\hat{h} = \hat{H}(X_t) - \hat{I}$ revealed that 99.6\% of the entropy-rate offset originated in $\hat{H}$, not in $\hat{I}$.

\subsubsection{Root cause}
The ANN library's k-d tree search (\texttt{kd\_search.cpp}, line~234) contains the guard
\begin{lstlisting}[language=C,numbers=none]
if (d >= ANNkdDim && (ANN_ALLOW_SELF_MATCH || dist != 0))
    ANNkdPointMK->insert(dist, bkt[i]);
\end{lstlisting}
With \texttt{ANN\_ALLOW\_SELF\_MATCH = false}, the condition \texttt{dist != 0} excludes \emph{all} points at zero distance from the result set---not only the query point.
For integer-valued data with many duplicates, this guarantees $\varepsilon_i > 0$ for every query, so $N_{\mathrm{v}} = \Neff$ in Eq.~\eqref{eq:kl}.

In contrast, Python's \texttt{scipy.spatial.KDTree} returns exact-zero distances for duplicate points.
The original Python code excluded these as ``errors,'' reducing $N_{\mathrm{v}}$ and shifting the $\psi(N_{\mathrm{v}})$ term.
With $k = 5$ and typical data: $\psi(123) - \psi(200) \approx -0.49$, closely matching the observed offset.

\subsubsection{Correction}
We modified \texttt{\_entropy\_knn} to count the maximum number of co-located points via \texttt{np.unique}, request $k + d_{\max}$ neighbors (where $d_{\max}$ is the maximum duplicate count), and for each query point select the $k$-th neighbor with non-zero distance.
This is implemented via vectorized NumPy operations (Algorithm~\ref{alg:skipdup}).

\begin{algorithm}[t]
\caption{Zero-distance--aware k-NN entropy estimation}\label{alg:skipdup}
\begin{algorithmic}[1]
\Require Data $X \in \mathbb{R}^{d \times N}$, neighbors $k$
\State $d_{\max} \gets \max(\text{duplicate counts via \textsc{Unique}})$
\State $k_{\mathrm{ext}} \gets \min(N,\; k + d_{\max})$
\State $D \gets \textsc{KDTree}(X).\textsc{Query}(X, k_{\mathrm{ext}}, p = \infty)$ \Comment{$N \times k_{\mathrm{ext}}$ distance matrix}
\State $n_{\mathrm{zero},i} \gets \sum_{j} \mathbb{1}[D_{ij} = 0]$ \Comment{zeros per row (incl.\ self)}
\State $\varepsilon_i \gets D_{i,\, n_{\mathrm{zero},i} + k - 1}$ \Comment{$k$-th non-zero distance}
\State $N_{\mathrm{v}} \gets \sum_i \mathbb{1}[\varepsilon_i > 0]$
\State \Return $d \cdot \overline{\log(2\varepsilon)}_{i:\varepsilon_i > 0} + \psi(N_{\mathrm{v}}) - \psi(k)$
\end{algorithmic}
\end{algorithm}

After correction, $H(X_t)$ matches to within numerical precision ($R^2 = 1.000$, mean absolute difference $< 10^{-12}$).


\subsection{Discrepancy 2: $L_2$ vs.\ $L_\infty$ norm}\label{sec:bug2}

\subsubsection{Observation}
Shannon entropy $H$ at embedding dimension $m = 2$ (a two-dimensional quantity) showed a consistent offset of \SI{0.25}{nats} across all time scales, despite excellent correlation ($R^2 = 0.9996$).

\subsubsection{Root cause}
The ANN library is compiled with \texttt{ANN\_METRIC = LINF} (\texttt{ANN.h}, lines~382--387), using the Chebyshev ($L_\infty$) norm for all distance computations.
Python's \texttt{scipy.spatial.KDTree.query} defaults to $L_2$ (Euclidean).
In one dimension $L_2 \equiv L_\infty$, so the marginal entropy $H(X_t)$ was unaffected.
However, for $d = 2$ the norms differ: $\|x\|_\infty \leq \|x\|_2 \leq \sqrt{d}\,\|x\|_\infty$, producing systematically larger distance estimates under $L_2$ and thus a positive entropy bias.

\subsubsection{Correction}
We added \texttt{p=np.inf} to the \texttt{tree.query()} call in \texttt{\_entropy\_knn}, selecting the Chebyshev norm.
After correction, Shannon entropy matches to machine precision across all scales ($R^2 = 1.0000$, Table~\ref{tab:results}).


\subsection{Discrepancy 3: Embedding misalignment in entropy rate}\label{sec:bug3}

\subsubsection{Observation}
After correcting Discrepancies~1 and~2, the entropy rate at $\tau = 1$ matched well ($R^2 = 0.9989$), but at $\tau = 2$ and $\tau = 4$ significant offsets persisted (mean absolute differences of 0.24 and 0.19, $R^2 \approx 0.85$; Table~\ref{tab:before_after}).

\subsubsection{Root cause}
The Python \texttt{compute\_entropy\_rate} (method~2) computed $\hat{H}(X_t)$ and $\hat{I}(X_t; X_{\mathrm{past}})$ via separate function calls.
Two problems emerged at $\tau > 1$:

\begin{enumerate}
  \item \textbf{Different data for $H$ and $I$.}
  The call \texttt{compute\_entropy(x, n\_embed=1, stride=$\tau$)} selected the first $\Neff$ \emph{consecutive} points of the raw signal (since \texttt{pts\_offset} $= \tau \cdot (1-1) = 0$ for \texttt{n\_embed}${}=1$).
  Meanwhile, the MI computation used points from a time-delay embedding that starts at offset $\tau \cdot m$.
  The subtraction $\hat{h} = \hat{H} - \hat{I}$ is only meaningful when both terms operate on the same data points.

  \item \textbf{Misaligned current/past pairing.}
  The past vectors were generated by \texttt{\_embed\_data(x, m, $\tau$)}, which steps by $\tau$ between consecutive embedded vectors.
  However, the ``current'' values \texttt{x\_curr = x[:, offset:]} stepped by~1.
  At $\tau = 2$ this caused pairing errors: the $i$-th current value $x_{4+i}$ was paired with the past vector for time $2i + 4$, matching only for $i = 0$.
\end{enumerate}

The C implementation avoids both problems by constructing a single unified embedding of $(m+1)$ dimensions via \texttt{Theiler\_embed}, from which the current values (first row) and past values (remaining rows) are extracted.
All $\Neff$ embedded vectors step by $\Delta_T = 1$ (the Theiler window), with temporal lags of~$\tau$ between embedding dimensions.

\subsubsection{Correction}
We rewrote \texttt{compute\_entropy\_rate} method~2 to construct the unified embedding directly:
\begin{enumerate}
  \item Compute $t_i = \tau \cdot m + i$ for $i = 0, \ldots, \Neff - 1$ (consecutive time indices).
  \item Extract $X_{\mathrm{curr}} = \{x_{t_i}\}$ and $X_{\mathrm{past}} = \{(x_{t_i - \tau}, \ldots, x_{t_i - m\tau})\}$.
  \item Compute $\hat{H}$ on $X_{\mathrm{curr}}$ and $\hat{I}$ on $(X_{\mathrm{curr}}, X_{\mathrm{past}})$.
\end{enumerate}

After correction, entropy-rate estimates agree to within 0.06\% at all scales (Table~\ref{tab:results}).


\subsection{Final validation results}\label{sec:final}

\begin{table}[t]
\centering
\caption{Comparison of Python and C implementations after all corrections. Shannon entropy $H(m{=}2)$ and entropy rate $h(m{=}2)$ across 250 valid segments at three time scales.}\label{tab:results}
\begin{tabular}{@{}llcccc@{}}
\toprule
Metric & Scale & Mean abs.\ diff.\ & Mean rel.\ \% & Max abs.\ diff.\ & $R^2$ \\
\midrule
\multirow{3}{*}{Shannon $H$} & $\tau = 1$ & $< 10^{-12}$ & $0.00$ & $< 10^{-12}$ & $1.0000$ \\
                              & $\tau = 2$ & $< 10^{-12}$ & $0.00$ & $< 10^{-12}$ & $1.0000$ \\
                              & $\tau = 4$ & $< 10^{-12}$ & $0.00$ & $< 10^{-12}$ & $1.0000$ \\
\midrule
\multirow{3}{*}{Entropy rate $h$} & $\tau = 1$ & $0.0029$ & $0.06$ & $0.1084$ & $0.9997$ \\
                                  & $\tau = 2$ & $0.0014$ & $0.03$ & $0.0371$ & $0.9999$ \\
                                  & $\tau = 4$ & $0.0007$ & $0.01$ & $0.0242$ & $0.9999$ \\
\bottomrule
\end{tabular}
\end{table}

\begin{table}[t]
\centering
\caption{Progression of entropy-rate agreement as each discrepancy is corrected (values at $\tau = 1$).}\label{tab:before_after}
\begin{tabular}{@{}lccc@{}}
\toprule
Stage & Mean abs.\ diff.\ & $R^2$ & Primary cause \\
\midrule
Before any corrections & $0.4302$ & $0.836$ & All three \\
After zero-distance fix & $0.0113$ & $0.9989$ & Norm + alignment \\
After $L_\infty$ norm fix & $0.0029$ & $0.9997$ & Alignment (minor at $\tau{=}1$) \\
After embedding alignment & $0.0029$ & $0.9997$ & Residual floating-point \\
\bottomrule
\end{tabular}
\end{table}

Shannon entropy $H(m{=}2)$ now matches to machine precision at all tested time scales (Table~\ref{tab:results}).
The residual entropy-rate difference of ${\sim}0.003$ at $\tau = 1$ reflects minor differences in the marginal-count step of the KSG estimator (Algorithm~1), where the C code uses an iterative point-counting approach inside the ANN tree structure, while Python uses \texttt{scipy}'s \texttt{query\_ball\_point}.
At $\tau = 2$ and $\tau = 4$ the agreement is even tighter (mean difference $< 0.002$), likely because the coarser embedding reduces sensitivity to counting-boundary effects.


\section{Discussion}\label{sec:discussion}

\subsection{Practical implications}

The three discrepancies identified in this work are not bugs in the mathematical sense---both the C and the original Python code implement valid entropy estimators---but they prevent numerical reproducibility across codebases.
We summarize the lessons for implementers:

\begin{enumerate}
  \item \textbf{Self-match exclusion vs.\ duplicate exclusion.}
  The ANN library's \texttt{dist != 0} guard was designed to exclude the query point from its own neighbor list.
  However, it also excludes all co-located (duplicate) points.
  Any reimplementation that wishes to match ANN's behavior must replicate this side effect explicitly.
  For continuous-valued data without ties, the distinction is immaterial; for quantized data (integer RR intervals, discrete event times), it dominates the entropy estimate.

  \item \textbf{Distance metric consistency.}
  The KL estimator's $\log(2\varepsilon)$ term assumes the $L_\infty$ ball volume $V_d = (2\varepsilon)^d$.
  When $L_2$ distances are used instead, the estimator implicitly uses $L_2$-ball volumes $V_d = \pi^{d/2} \varepsilon^d / \Gamma(d/2 + 1)$, introducing a dimension-dependent bias.
  The $L_\infty$ norm is both the correct choice for the KL estimator as commonly stated and the one used by the ANN library.

  \item \textbf{Unified embedding for entropy-rate decomposition.}
  The identity $h = H(X_t) - I(X_t; X_{\mathrm{past}})$ holds in expectation but is sensitive to finite-sample alignment.
  Computing $H$ and $I$ on different subsets of the data introduces a stochastic bias that does not vanish with $N$, because the digamma terms depend on the specific point set.
  A unified embedding that feeds both estimators from the same data matrix eliminates this source of error.
\end{enumerate}

\subsection{Remaining discrepancies}

The Sample Entropy (SampEn) estimates show weak agreement (Spearman $\rho = 0.20$) between the two implementations.
This reflects genuinely different algorithmic conventions in the template-matching step (inclusive vs.\ exclusive boundary, self-match handling) rather than a bug, and is not pursued further here.

\subsection{Performance}

The pure Python implementation runs approximately 12$\times$ slower than the C/Cython version on this dataset (8.3\,s vs.\ 0.7\,s for 251 segments).
This is acceptable for research and validation purposes but may require the optimized \texttt{core\_optimized} module (using \texttt{scipy.spatial.cKDTree} and parallel workers) for production workloads.


\section{Conclusion}\label{sec:conclusion}

We have validated a pure Python implementation of k-NN entropy and entropy-rate estimators against a reference C/ANN implementation, achieving numerical equivalence ($R^2 > 0.999$) on real-world physiological time-series data.
The three implementation-level discrepancies identified---zero-distance handling, distance-metric mismatch, and embedding misalignment---are general pitfalls that apply to any reimplementation of k-NN information-theoretic estimators.
The corrected Python codebase provides a readable, dependency-minimal reference implementation suitable for research, teaching, and cross-validation of compiled libraries.

All code and validation scripts are available at \url{https://github.com/martinfrasch/entropy} (branch \texttt{claude/refactor-python-efficiency-UGyKA}).


\section*{Acknowledgments}

The systematic debugging and validation was performed with assistance from Claude~Code (Anthropic), which identified the ANN library's zero-distance exclusion mechanism as the primary root cause.


\bibliographystyle{abbrvnat}
\bibliography{references}

\end{document}
